% =============================================================================
% Conclusion -- three-claim framing
% Drop-in include: % =============================================================================
% Conclusion -- three-claim framing
% Drop-in include: % =============================================================================
% Conclusion -- three-claim framing
% Drop-in include: % =============================================================================
% Conclusion -- three-claim framing
% Drop-in include: \input{sections/conclusion}
% =============================================================================
\section{Conclusion}
\label{sec:conclusion}

This study does not show that GRPO generally improves reasoning, and we are
careful not to claim that it does.  The training runner is a group-relative
policy-gradient variant inspired by GRPO rather than a faithful implementation
of the canonical objective (missing PPO clip, reference-KL, and completion-only
mask; a P1-A diagnostic measures $\sim$76\% of the full-sequence loss magnitude
to be prompt-token gradient), so algorithm-level conclusions apply to this
runner and not to GRPO-the-algorithm.  Three of our reward environments
(HumanEval, tool-use, and MATH) measure sandbox-restricted execution, JSON
schema compliance, and boxed-answer formatting respectively, not end-to-end
capability.  The MATH-500 and HumanEval prompt pools in the training loop are
loaded from the \texttt{test} splits and function as reward-environment
probes, not as generalization tests.  The only capability signal evaluated on
a genuinely held-out slice is GSM8K, where the lift from $82.0\%$ to
$83.3\%$ is not significant ($p=0.26$).  The defensible conclusion is that
this critic-free group-relative RL loop behaves as a conditional reward-contrast
amplifier: it can refine behavior when the current policy samples both
successes and failures under a verifier that rewards them, and it cannot
reliably bootstrap competence when sampled groups are flat or the verifier
rejects legitimate completions.

The paper's main contribution is therefore diagnostic.  ZVF and GU should be
logged early and interpreted jointly with reward level.  High ZVF with low
reward is a cold-start alarm; high ZVF with high reward is saturation; low ZVF
means rollout compute is still producing contrast.  In our validation set,
the first-five-step rule catches the two collapsed runs without false
positives, but the low number of failures and weak continuous-ranking result
mean ZVF/GU should be treated as triage diagnostics, not performance
predictors.

The mechanism behind that diagnostic is mixed-group probability:
\[
P(\text{usable}) \approx \frac{1}{N}\sum_x \left[1 - (1-p_x)^G - p_x^G\right].
\]
This formula is the cleanest way to state the boundary condition.  The
expected-one-success rule $p_0>1/G$ is only a heuristic.  Future work should
test the three direct interventions implied by the equation: raise $p_0$ with
SFT or easier curricula, raise $G$ where rollout cost allows, and densify the
reward so equivalent partial progress is not collapsed into all-zero groups.

Finally, algorithm labels should be reported as stack-conditioned results.
The Qwen PPO/GRPO row is artifact-sensitive, and the Llama row favors PPO;
neither supports a universal PPO/GRPO ordering.  A publishable comparison
needs the full treatment definition: model family, backend, reward parser,
sampler, optimizer, KL/reference handling, LoRA configuration, checkpoint
selection, and held-out evaluation.  Without those controls, ``PPO versus
GRPO'' is not an interpretable scientific claim.

% --------------------------------------------------------------------------=
% Claims we do not make (from no-go list)
% This section explicitly enumerates claims we explicitly avoid making.
% --------------------------------------------------------------------------=

\section*{Claims We Do Not Make}

We do not claim that our runner is a faithful implementation of canonical GRPO
or that the results constitute an algorithm leaderboard: the experiments measure
a particular stack-conditioned, group-relative policy-gradient loop with specific
choices about masking, KL/reference handling, sampling, LoRA, optimization,
backend, and evaluation.  We do not claim that GRPO broadly improves reasoning,
alignment, coding, tool use, browser control, or general LLM capability, nor
that training reward or post-selected high-reward checkpoints establish causal
held-out gains; the only clean held-out capability check is GSM8K, where the
82.0\% to 83.3\% lift is not statistically significant.  We do not treat
MATH-500 or HumanEval reward-environment probes, boxed-answer rewards,
restricted HumanEval sandboxes, JSON-schema tool rewards, or BrowserGym smoke
tests as evidence of end-to-end generalization, execution success, or agentic
competence.  We do not claim universal PPO/GRPO/DPO, Tinker/TRL, dense/MoE,
base/instruct, frontier-scale, group-size, or scaling-law rankings, because
many comparisons are single-seed, partial, post-selected, artifact-sensitive,
closed-source, or confounded by stack details.  We do not claim that ZVF/GU
are causal or calibrated predictors of final performance, only that they are
useful triage diagnostics for cold-start collapse, saturation, and within-group
reward contrast.  We do not claim that verifiable rewards eliminate reward
hacking, that reward-stability proxies are direct KL measurements, that short
30-step runs characterize asymptotic behavior, or that these verifiable-reward
LoRA experiments extrapolate to human-preference RLHF, safety, truthfulness,
real tool execution, or open-ended deployment.

% =============================================================================
\section{Conclusion}
\label{sec:conclusion}

This study does not show that GRPO generally improves reasoning, and we are
careful not to claim that it does.  The training runner is a group-relative
policy-gradient variant inspired by GRPO rather than a faithful implementation
of the canonical objective (missing PPO clip, reference-KL, and completion-only
mask; a P1-A diagnostic measures $\sim$76\% of the full-sequence loss magnitude
to be prompt-token gradient), so algorithm-level conclusions apply to this
runner and not to GRPO-the-algorithm.  Three of our reward environments
(HumanEval, tool-use, and MATH) measure sandbox-restricted execution, JSON
schema compliance, and boxed-answer formatting respectively, not end-to-end
capability.  The MATH-500 and HumanEval prompt pools in the training loop are
loaded from the \texttt{test} splits and function as reward-environment
probes, not as generalization tests.  The only capability signal evaluated on
a genuinely held-out slice is GSM8K, where the lift from $82.0\%$ to
$83.3\%$ is not significant ($p=0.26$).  The defensible conclusion is that
this critic-free group-relative RL loop behaves as a conditional reward-contrast
amplifier: it can refine behavior when the current policy samples both
successes and failures under a verifier that rewards them, and it cannot
reliably bootstrap competence when sampled groups are flat or the verifier
rejects legitimate completions.

The paper's main contribution is therefore diagnostic.  ZVF and GU should be
logged early and interpreted jointly with reward level.  High ZVF with low
reward is a cold-start alarm; high ZVF with high reward is saturation; low ZVF
means rollout compute is still producing contrast.  In our validation set,
the first-five-step rule catches the two collapsed runs without false
positives, but the low number of failures and weak continuous-ranking result
mean ZVF/GU should be treated as triage diagnostics, not performance
predictors.

The mechanism behind that diagnostic is mixed-group probability:
\[
P(\text{usable}) \approx \frac{1}{N}\sum_x \left[1 - (1-p_x)^G - p_x^G\right].
\]
This formula is the cleanest way to state the boundary condition.  The
expected-one-success rule $p_0>1/G$ is only a heuristic.  Future work should
test the three direct interventions implied by the equation: raise $p_0$ with
SFT or easier curricula, raise $G$ where rollout cost allows, and densify the
reward so equivalent partial progress is not collapsed into all-zero groups.

Finally, algorithm labels should be reported as stack-conditioned results.
The Qwen PPO/GRPO row is artifact-sensitive, and the Llama row favors PPO;
neither supports a universal PPO/GRPO ordering.  A publishable comparison
needs the full treatment definition: model family, backend, reward parser,
sampler, optimizer, KL/reference handling, LoRA configuration, checkpoint
selection, and held-out evaluation.  Without those controls, ``PPO versus
GRPO'' is not an interpretable scientific claim.

% --------------------------------------------------------------------------=
% Claims we do not make (from no-go list)
% This section explicitly enumerates claims we explicitly avoid making.
% --------------------------------------------------------------------------=

\section*{Claims We Do Not Make}

We do not claim that our runner is a faithful implementation of canonical GRPO
or that the results constitute an algorithm leaderboard: the experiments measure
a particular stack-conditioned, group-relative policy-gradient loop with specific
choices about masking, KL/reference handling, sampling, LoRA, optimization,
backend, and evaluation.  We do not claim that GRPO broadly improves reasoning,
alignment, coding, tool use, browser control, or general LLM capability, nor
that training reward or post-selected high-reward checkpoints establish causal
held-out gains; the only clean held-out capability check is GSM8K, where the
82.0\% to 83.3\% lift is not statistically significant.  We do not treat
MATH-500 or HumanEval reward-environment probes, boxed-answer rewards,
restricted HumanEval sandboxes, JSON-schema tool rewards, or BrowserGym smoke
tests as evidence of end-to-end generalization, execution success, or agentic
competence.  We do not claim universal PPO/GRPO/DPO, Tinker/TRL, dense/MoE,
base/instruct, frontier-scale, group-size, or scaling-law rankings, because
many comparisons are single-seed, partial, post-selected, artifact-sensitive,
closed-source, or confounded by stack details.  We do not claim that ZVF/GU
are causal or calibrated predictors of final performance, only that they are
useful triage diagnostics for cold-start collapse, saturation, and within-group
reward contrast.  We do not claim that verifiable rewards eliminate reward
hacking, that reward-stability proxies are direct KL measurements, that short
30-step runs characterize asymptotic behavior, or that these verifiable-reward
LoRA experiments extrapolate to human-preference RLHF, safety, truthfulness,
real tool execution, or open-ended deployment.

% =============================================================================
\section{Conclusion}
\label{sec:conclusion}

This study does not show that GRPO generally improves reasoning, and we are
careful not to claim that it does.  The training runner is a group-relative
policy-gradient variant inspired by GRPO rather than a faithful implementation
of the canonical objective (missing PPO clip, reference-KL, and completion-only
mask; a P1-A diagnostic measures $\sim$76\% of the full-sequence loss magnitude
to be prompt-token gradient), so algorithm-level conclusions apply to this
runner and not to GRPO-the-algorithm.  Three of our reward environments
(HumanEval, tool-use, and MATH) measure sandbox-restricted execution, JSON
schema compliance, and boxed-answer formatting respectively, not end-to-end
capability.  The MATH-500 and HumanEval prompt pools in the training loop are
loaded from the \texttt{test} splits and function as reward-environment
probes, not as generalization tests.  The only capability signal evaluated on
a genuinely held-out slice is GSM8K, where the lift from $82.0\%$ to
$83.3\%$ is not significant ($p=0.26$).  The defensible conclusion is that
this critic-free group-relative RL loop behaves as a conditional reward-contrast
amplifier: it can refine behavior when the current policy samples both
successes and failures under a verifier that rewards them, and it cannot
reliably bootstrap competence when sampled groups are flat or the verifier
rejects legitimate completions.

The paper's main contribution is therefore diagnostic.  ZVF and GU should be
logged early and interpreted jointly with reward level.  High ZVF with low
reward is a cold-start alarm; high ZVF with high reward is saturation; low ZVF
means rollout compute is still producing contrast.  In our validation set,
the first-five-step rule catches the two collapsed runs without false
positives, but the low number of failures and weak continuous-ranking result
mean ZVF/GU should be treated as triage diagnostics, not performance
predictors.

The mechanism behind that diagnostic is mixed-group probability:
\[
P(\text{usable}) \approx \frac{1}{N}\sum_x \left[1 - (1-p_x)^G - p_x^G\right].
\]
This formula is the cleanest way to state the boundary condition.  The
expected-one-success rule $p_0>1/G$ is only a heuristic.  Future work should
test the three direct interventions implied by the equation: raise $p_0$ with
SFT or easier curricula, raise $G$ where rollout cost allows, and densify the
reward so equivalent partial progress is not collapsed into all-zero groups.

Finally, algorithm labels should be reported as stack-conditioned results.
The Qwen PPO/GRPO row is artifact-sensitive, and the Llama row favors PPO;
neither supports a universal PPO/GRPO ordering.  A publishable comparison
needs the full treatment definition: model family, backend, reward parser,
sampler, optimizer, KL/reference handling, LoRA configuration, checkpoint
selection, and held-out evaluation.  Without those controls, ``PPO versus
GRPO'' is not an interpretable scientific claim.

% --------------------------------------------------------------------------=
% Claims we do not make (from no-go list)
% This section explicitly enumerates claims we explicitly avoid making.
% --------------------------------------------------------------------------=

\section*{Claims We Do Not Make}

We do not claim that our runner is a faithful implementation of canonical GRPO
or that the results constitute an algorithm leaderboard: the experiments measure
a particular stack-conditioned, group-relative policy-gradient loop with specific
choices about masking, KL/reference handling, sampling, LoRA, optimization,
backend, and evaluation.  We do not claim that GRPO broadly improves reasoning,
alignment, coding, tool use, browser control, or general LLM capability, nor
that training reward or post-selected high-reward checkpoints establish causal
held-out gains; the only clean held-out capability check is GSM8K, where the
82.0\% to 83.3\% lift is not statistically significant.  We do not treat
MATH-500 or HumanEval reward-environment probes, boxed-answer rewards,
restricted HumanEval sandboxes, JSON-schema tool rewards, or BrowserGym smoke
tests as evidence of end-to-end generalization, execution success, or agentic
competence.  We do not claim universal PPO/GRPO/DPO, Tinker/TRL, dense/MoE,
base/instruct, frontier-scale, group-size, or scaling-law rankings, because
many comparisons are single-seed, partial, post-selected, artifact-sensitive,
closed-source, or confounded by stack details.  We do not claim that ZVF/GU
are causal or calibrated predictors of final performance, only that they are
useful triage diagnostics for cold-start collapse, saturation, and within-group
reward contrast.  We do not claim that verifiable rewards eliminate reward
hacking, that reward-stability proxies are direct KL measurements, that short
30-step runs characterize asymptotic behavior, or that these verifiable-reward
LoRA experiments extrapolate to human-preference RLHF, safety, truthfulness,
real tool execution, or open-ended deployment.

% =============================================================================
\section{Conclusion}
\label{sec:conclusion}

This study does not show that GRPO generally improves reasoning, and we are
careful not to claim that it does.  The training runner is a group-relative
policy-gradient variant inspired by GRPO rather than a faithful implementation
of the canonical objective (missing PPO clip, reference-KL, and completion-only
mask; a P1-A diagnostic measures $\sim$76\% of the full-sequence loss magnitude
to be prompt-token gradient), so algorithm-level conclusions apply to this
runner and not to GRPO-the-algorithm.  Three of our reward environments
(HumanEval, tool-use, and MATH) measure sandbox-restricted execution, JSON
schema compliance, and boxed-answer formatting respectively, not end-to-end
capability.  The MATH-500 and HumanEval prompt pools in the training loop are
loaded from the \texttt{test} splits and function as reward-environment
probes, not as generalization tests.  The only capability signal evaluated on
a genuinely held-out slice is GSM8K, where the lift from $82.0\%$ to
$83.3\%$ is not significant ($p=0.26$).  The defensible conclusion is that
this critic-free group-relative RL loop behaves as a conditional reward-contrast
amplifier: it can refine behavior when the current policy samples both
successes and failures under a verifier that rewards them, and it cannot
reliably bootstrap competence when sampled groups are flat or the verifier
rejects legitimate completions.

The paper's main contribution is therefore diagnostic.  ZVF and GU should be
logged early and interpreted jointly with reward level.  High ZVF with low
reward is a cold-start alarm; high ZVF with high reward is saturation; low ZVF
means rollout compute is still producing contrast.  In our validation set,
the first-five-step rule catches the two collapsed runs without false
positives, but the low number of failures and weak continuous-ranking result
mean ZVF/GU should be treated as triage diagnostics, not performance
predictors.

The mechanism behind that diagnostic is mixed-group probability:
\[
P(\text{usable}) \approx \frac{1}{N}\sum_x \left[1 - (1-p_x)^G - p_x^G\right].
\]
This formula is the cleanest way to state the boundary condition.  The
expected-one-success rule $p_0>1/G$ is only a heuristic.  Future work should
test the three direct interventions implied by the equation: raise $p_0$ with
SFT or easier curricula, raise $G$ where rollout cost allows, and densify the
reward so equivalent partial progress is not collapsed into all-zero groups.

Finally, algorithm labels should be reported as stack-conditioned results.
The Qwen PPO/GRPO row is artifact-sensitive, and the Llama row favors PPO;
neither supports a universal PPO/GRPO ordering.  A publishable comparison
needs the full treatment definition: model family, backend, reward parser,
sampler, optimizer, KL/reference handling, LoRA configuration, checkpoint
selection, and held-out evaluation.  Without those controls, ``PPO versus
GRPO'' is not an interpretable scientific claim.

% --------------------------------------------------------------------------=
% Claims we do not make (from no-go list)
% This section explicitly enumerates claims we explicitly avoid making.
% --------------------------------------------------------------------------=

\section*{Claims We Do Not Make}

We do not claim that our runner is a faithful implementation of canonical GRPO
or that the results constitute an algorithm leaderboard: the experiments measure
a particular stack-conditioned, group-relative policy-gradient loop with specific
choices about masking, KL/reference handling, sampling, LoRA, optimization,
backend, and evaluation.  We do not claim that GRPO broadly improves reasoning,
alignment, coding, tool use, browser control, or general LLM capability, nor
that training reward or post-selected high-reward checkpoints establish causal
held-out gains; the only clean held-out capability check is GSM8K, where the
82.0\% to 83.3\% lift is not statistically significant.  We do not treat
MATH-500 or HumanEval reward-environment probes, boxed-answer rewards,
restricted HumanEval sandboxes, JSON-schema tool rewards, or BrowserGym smoke
tests as evidence of end-to-end generalization, execution success, or agentic
competence.  We do not claim universal PPO/GRPO/DPO, Tinker/TRL, dense/MoE,
base/instruct, frontier-scale, group-size, or scaling-law rankings, because
many comparisons are single-seed, partial, post-selected, artifact-sensitive,
closed-source, or confounded by stack details.  We do not claim that ZVF/GU
are causal or calibrated predictors of final performance, only that they are
useful triage diagnostics for cold-start collapse, saturation, and within-group
reward contrast.  We do not claim that verifiable rewards eliminate reward
hacking, that reward-stability proxies are direct KL measurements, that short
30-step runs characterize asymptotic behavior, or that these verifiable-reward
LoRA experiments extrapolate to human-preference RLHF, safety, truthfulness,
real tool execution, or open-ended deployment.
